\section{The rate for adjacent coefficient ratios}
\label{sec:quantitative-coefficients}

This appendix proves Theorem~\ref{thm:quantitative-coefficient}, which
contains the rate in Theorem~\ref{thm:coefficient}. We quantify the three
degree ranges treated in Section~\ref{uniform-one-boundary-coefficient-ratios}.
The logarithmic comparison and concavity give the rate away from the
lowest degrees, including bounded positive codimension at the top.
Bracket difference estimates supply the remaining range. Balancing the
bounds gives $O(g^{-1/8})$ uniformly; we do not claim that this exponent
is optimal. The coefficients and the exact density are those of
Theorem~\ref{thm:coefficient}.

\begin{theorem}[Quantitative adjacent ratios]\label{thm:quantitative-coefficient}
There is a constant $C$ such that, for all sufficiently large $g$,
\begin{equation}
 \sup_{0\le k<3g-2}
 \left|\frac{a_{g,k}/a_{g,k+1}}{4\pi^2/t_{g,k}^2}-1\right|
 \le Cg^{-1/8}.
 \label{eq:QC1}
\end{equation}
For every fixed $\varepsilon\in(0,3)$ there is $C_\varepsilon$ such that
\begin{equation}
 \sup_{\varepsilon g\le k<3g-2}
 \left|\frac{a_{g,k}/a_{g,k+1}}{4\pi^2/t_{g,k}^2}-1\right|
 \le C_\varepsilon g^{-1/4}.
 \label{eq:QC2}
\end{equation}
In the second supremum $k$ is an integer. Both bounds include bounded
positive codimension from the top coefficient.
\end{theorem}

\subsection{The logarithmic model up to its endpoints}

We first make the logarithmic comparison of
Section~\ref{uniform-one-boundary-coefficient-ratios} quantitative, up to
both ends of the degree range.

Use $H,\widetilde A_m,\mu$ from \eqref{eq:C4} and \eqref{eq:C10}.
Let $\zeta(u)$ be the inverse of $\mu(z)$ on $0\le u\le3/2$,
with $\zeta(0)=0$. Since $\mu(z)=z/5+O(z^2)$ at zero and
$\mu'(z)>0$ throughout $[0,\pi^2/2]$, the function
\[
 \chi(u)=u/\zeta(u),\qquad \chi(0)=1/5,
\]
is positive and analytic in a neighborhood of the closed real interval.
Write $f$ for the function in \eqref{eq:C14}, extended continuously to
$[0,3]$, and set
\begin{equation}
 \mathfrak q(\alpha)=4\alpha\chi((3-\alpha)/2),\qquad
 f'(\alpha)=-\log\mathfrak q(\alpha),\qquad
 |f''(\alpha)|\le C(1+\alpha^{-1})\quad(0<\alpha\le3).
 \label{eq:QC3}
\end{equation}
These follow from \eqref{eq:C15}; in particular
$\mathfrak q(3)=12/5$ and $f(3)=0$.

We need the endpoint-uniform version of the logarithmic saddle estimate
\eqref{eq:C12}.

\begin{lemma}[Endpoint-uniform saddle estimate]\label{lem:uniform-saddle}
For integers $0\le m\le3h/2$, let $z=\zeta(m/h)$, interpreting the
$m=0$ expressions by continuity. Then, uniformly in $m$,
\begin{equation}
 \log \widetilde A_m(h)=-h\log H(z)-m\log z+O(\log(h+1)).
 \label{eq:QC4}
\end{equation}
\end{lemma}

\begin{proof}
Tilt $H^{-1}$ at $z$ and let
$S_h$ be the sum of $h$ independent variables with probability
generating function $H(z)/H(zw)$. Exactly,
\[
 \widetilde A_m(h)=H(z)^{-h}z^{-m}\Pr(S_h=m).
\]
There is $c>0$ such that $\Pr(S_h=m)\ge c/\sqrt{m+1}$ uniformly.
We give the details at the degenerate endpoint $m/h\to0$.
When $m\to\infty$ and $m/h\to0$, one has $z\asymp m/h$,
$\operatorname{Var}(S_h)\asymp m$, and each fixed cumulant of the
sum is $O(m)$. The probabilities at zero and one of one summand have
product comparable to $z$. The Fourier bound \eqref{eq:TL40} therefore
gives $\exp(-cm\theta^2)$ on $[-\pi,\pi]$, and the centered Taylor
expansion gives the local Gaussian limit at $m$, with central mass
asymptotic to $(2\pi\operatorname{Var}(S_h))^{-1/2}$.
When $m/h$ stays away from zero, the same argument uses a compact
tilted family with variance comparable to $h$ and positive zero and
one probabilities. If $m$ stays bounded, pass to a fixed integer $m$;
then $z=5m/h+O_m(h^{-2})$ and
$\widetilde A_m(h)=(h/5)^m/m!\,(1+O_m(h^{-1}))$, so the central mass tends to
$\mathrm{e}^{-m}m^m/m!>0$. At $m=0$ it is one.
These cases and a subsequence contradiction prove the asserted uniform
lower bound. The upper bound is one, proving \eqref{eq:QC4}.
\end{proof}

Put
\[
 G_g(k)=\log\frac{J_k}{J_Ng^{2(N-k)}},\qquad N=3g-2.
\]

\begin{lemma}[Logarithmic comparison up to the endpoints]\label{lem:log-comparison}
There is a constant $C$ such that
\begin{equation}
 \max_{0\le k\le N}|G_g(k)-g f(k/g)|\le C\sqrt g.
 \label{eq:QC5}
\end{equation}
\end{lemma}

\begin{proof}
This follows from the model \eqref{eq:C4}, the comparison \eqref{eq:C5} and
Lemma~\ref{lem:uniform-saddle}. The passage to $f$ is uniform, even at
$m=0$. To see this, use the
exact model for the intersection number,
$3^h\Gamma(h)4^{-g}m!\widetilde A_m(h)/(2k+1)!!$ from Section~\ref{uniform-one-boundary-coefficient-ratios}
and apply Stirling's bound with logarithmic error to its factorials,
using $0\log0=0$. Replacing $h=2g-1$ by $2g$ and
$m=3g-2-k$ by $3g-k$ changes its saddle action by $O(\log(g+1))$:
the two partial derivatives of the action are $-\log H(z)$ and
$-\log z$, the latter having only an integrable logarithmic
singularity at $m=0$. The powers of $g$ cancel as in
\eqref{eq:C13}--\eqref{eq:C14}. Subtracting the top model value gives
$g f(k/g)+O(\log(g+1))$ uniformly. The two logarithmic errors between
the true and the model values are at most $C\sqrt g$ by \eqref{eq:C5}, including the top
coefficient, which proves \eqref{eq:QC5}.
\end{proof}

\subsection{Quantitative secants and the top endpoint}

\begin{proof}[Proof of Theorem~\ref{thm:quantitative-coefficient}]
The sequence $G_g$ is concave by \eqref{eq:C7}. Write
$q_k=J_k/J_{k+1}$. For an integer $\ell\ge1$ with
$\ell\le k/2$ and $k+\ell\le N$, its secants imply
\[
 -\frac{G_g(k)-G_g(k-\ell)}\ell
 \le\log(q_k/g^2)
 \le-\frac{G_g(k+\ell)-G_g(k)}\ell.
\]
Equations \eqref{eq:QC3} and \eqref{eq:QC5} yield
\begin{equation}
 \left|\log\frac{q_k}{g^2\mathfrak q(k/g)}\right|
 \le C\left(\frac{\sqrt g}\ell+\frac\ell k\right).
 \label{eq:QC6}
\end{equation}
If $k+\ell>N$, retain the lower secant and use
$q_k\le q_{N-1}$ for the upper bound. The exact top anchor
\eqref{eq:C18} gives $q_{N-1}/g^2=12/5+O(g^{-1})$.
Since $\log\mathfrak q$ is Lipschitz near three and
$3-k/g\le(\ell+2)/g$, this gives the same bound
\eqref{eq:QC6}, with a changed constant.

Take $\ell=\lfloor g^{5/8}\rfloor$. For $k\ge g^{3/4}$ it is
at most $k/2$ for large $g$, and \eqref{eq:QC6} is $O(g^{-1/8})$.
For $k\ge\varepsilon g$, instead take
$\ell=\lfloor g^{3/4}\rfloor$ to obtain $O_\varepsilon(g^{-1/4})$.

These estimates for the intersection numbers translate into the
statement of the theorem without losing the bounded-codimension range.
Indeed, put
$m=N-k$, $h=2g-1$, and $z_{g,k}=t_{g,k}^2/2$. By \eqref{eq:C11},
$\mu(z_{g,k})=m/h$. Therefore the value of $q_k=J_k/J_{k+1}$ that
corresponds, through \eqref{eq:C3}, to the proposed equivalent
$4\pi^2/t_{g,k}^2$ of $a_{g,k}/a_{g,k+1}$ is
\begin{equation}
 \frac{(2k+3)m}{z_{g,k}}
 =(2k+3)h\chi(m/h)
 =g^2\mathfrak q(k/g)\left(1+O(k^{-1}+g^{-1})\right).
 \label{eq:QC7}
\end{equation}
The last error is uniform for $k\ge1$: the two arguments of $\chi$
differ by $O(g^{-1})$, while $\chi$ is positive with bounded
logarithmic derivative on the closed interval. This calculation
retains $m$ exactly, including $m=1$. By \eqref{eq:C3}, the two
secant estimates prove \eqref{eq:QC1} for $k\ge g^{3/4}$ and all
of \eqref{eq:QC2}.

\subsection{The small-degree range}

It remains to consider $k\le g^{3/4}$. Put
$\rho_k=b_{g,k}/b_{g,k+1}\ge1$ and
$\omega_k=(2k+3)(N-k)$ as before. Equation~\eqref{eq:C25} gives
$\rho_k=1+O(g^{-1/2})$ for $k\le\delta_0\sqrt g$ with a fixed
small $\delta_0>0$. For the remaining indices we claim
\begin{equation}
 0\le\rho_k-1\le C\left(
       \frac{\log(g+1)}{\sqrt k}+\sqrt{\frac kg}+\frac1k\right)
 \quad(\delta_0\sqrt g\le k\le g^{3/4}).
 \label{eq:QC8}
\end{equation}
First $\rho_k<2$ uniformly in this range for large $g$. Otherwise
monotonicity of $q_j$ and $\omega_j/\omega_k\le1+\delta$ propagate a fixed
ratio gap over a block of length comparable to $k$. Its product
would force $\log(V_{g,1}/b_{g,j})\ge c k$ at an index $j\asymp k$,
contradicting \eqref{eq:C23}, whose upper bound is $o(k)$ uniformly
on this range.

Now put $\eta=\rho_k-1\in[0,1]$ and $\delta=\eta/6$.
For $k\le j\le\lfloor(1+\delta)k\rfloor$, the same monotonicity
and the exact factor $\omega_j$ give
$\rho_j\ge(1+\eta)/(1+\delta)\ge1+\eta/2$.
If $\eta k\ge12$, this block has at least $\eta k/12$ indices.
Summing logarithms over the block and using \eqref{eq:C23} at its
last index plus one gives
\[
 c\eta^2 k\le C\log^2(g+1)+Ck^2/g.
\]
All these indices lie below the fixed $c_0g$ range of \eqref{eq:C23}
for large $g$. If $\eta k<12$, then $\eta<12/k$ directly.
This proves \eqref{eq:QC8}. Its right side is $O(g^{-1/8})$ on
the displayed range. Thus $\rho_k=1+O(g^{-1/8})$ uniformly for
$0\le k\le g^{3/4}$.

Finally the saddle equation at $t=\pi$ gives
$\pi^2/t_{g,k}^2=1+O((k+1)/g)$ uniformly on this small-degree
range. Since $a_{g,k}/a_{g,k+1}=4\rho_k$, the desired
relative error is $O(g^{-1/8})$ there as well. Together with the
secant range this proves both assertions.
\end{proof}

The argument is a quantitative version of the one used for the leading
equivalent, based on the same logarithmic comparison and concavity. Liu--Xu's fixed-codimension expansion
has finer errors at each fixed codimension~\cite{LX}; the bounds here
hold with a single constant across all valid degrees.
